
\section{Results}
\subsection{RQ1: Smell Traceability}
\todo{The result box for RQ1 does not actually answer RQ1; it provides an observation rather than a concrete result.}
\ref{RQ1 Smell traceability} examines how many evolved smells can be traced to issues.
\autoref{tab:smell_linkage} presents the statistics on linkage for the analyzed projects, for both \Arcan and \ARCADE smells.

For \Arcan, we were not able to analyze all versions.
We encountered limitations with two versions of Shepard ("2024.08.05 Prepare quarkus migration", and " 2024.07.04 Monorepo ") and four versions of Shepard(0.5, 0.6, 0.7, 1.18.0).
\Arcan requires compiling the project and uses the Apache BCEL tool to read dependencies from jars and class files.
We were unable to compile these versions due to a missing dependency.
In contrast, \ARCADE operates only at the source code level, allowing analysis of all versions.

In general, \ARCADE smells are fewer and more traceable than \Arcan smells, %.\todo{proofread}
likely due to the different abstraction levels.
%It can be explained by their difference in abstraction level.
\Arcan operates at the Java class/package level, 
\ARCADE uses a clustering-based recovery approach for high-level components.
% \ARCADE uses a clustering-based architecture recovery approach to recover high-level components. % at a higher level of abstraction.
% On a higher level, smells are fewer but more traceable. 

Our traceability approach successfully linked many smells to their inducing issues, though not 100\%.
Coverage on Inkscape and StackGres is relatively poorer than on Shepard.
Some smells could not be traced due to development workflow practices and commits made outside of issue management.
This reflects the project characteristics shown in \autoref{tab:dataset}, Inkscape and StackGres contain far more commits inside the repository than in MR, showing that a lot of commits were outside the issue management system.

\begin{rqanswer}
\textbf{Answer to RQ1.}
Smells whose inducing changes occur within the issue management system can be linked to their inducing issues, but repositories may also contain smell-inducing changes outside issue management.
\end{rqanswer}



\subsection{RQ2: Prediction Performance}
\begin{table} 
    \centering
    \footnotesize
    \setlength{\tabcolsep}{1.43pt}
    \renewcommand{\arraystretch}{1.2}
    \caption{Results of the Top-2 Configuration of Each Classifier on the Smells Recovered With \ARCADE 
    (\acf{ffNN}, \acf{kNN}, \acf{LR}, \acf{NB}, \acf{RF}, \acf{SVM}, \acf{XGB}, and
    %ffNN = feedforward neural network,
    %kNN = k-nearest neighbors,
    %LR = logistic regression,
    %NB = Na\"ive Bayes,
    %RF = random forest,
    %SVM = support vector machine,
    %XGB = XGBoost
    Sentence-BERT (SB)).
    }
    \label{tab:rq2_arcade}
    \begin{tabularx}{\columnwidth}{@{}llcccccccccccc}
        \toprule
                                       &                          & \multicolumn{3}{c}{Inkscape} & \multicolumn{3}{c}{StackGres} & \multicolumn{3}{c}{Shepard} & \multicolumn{3}{c}{Avg.}                                                                                                           \\
        \cmidrule(lr){3-5} \cmidrule(lr){6-8} \cmidrule(lr){9-11} \cmidrule(lr){12-14}
        Class.                     & Emb.                & \multicolumn{1}{c}{P}                        & \multicolumn{1}{c}{R}                        & \multicolumn{1}{c}{F\textsubscript{1}}          & \multicolumn{1}{c}{P}                    & \multicolumn{1}{c}{R} & \multicolumn{1}{c}{F\textsubscript{1}} & \multicolumn{1}{c}{P} & \multicolumn{1}{c}{R} & \multicolumn{1}{c}{F\textsubscript{1}} & \multicolumn{1}{c}{P} & \multicolumn{1}{c}{R} & \multicolumn{1}{c}{F\textsubscript{1}} \\
        \midrule
\multirow{2}{*}{ffNN} & OpenAI  & .158 & .652 & .251 & .290 & .704 & .409 & .180 & .800 & .286 & .210 & .719 & .315 \\
 & SB\textsubscript{max}        & .150 & .591 & .229 & .275 & .704 & .392 & .169 & \textBF{.823} & .279 & .198 & .706 & .300 \\
\arrayrulecolor{kit-gray30}\hline\arrayrulecolor{black}
\multirow{2}{*}{kNN} & OpenAI   & .141 & \textBF{.694} & .235 & .274 & .751 & .401 & .178 & .816 & .292 & .198 & \textBF{.753} & .309 \\
 & SB\textsubscript{max}        & .134 & .589 & .218 & .260 & .692 & .378 & .206 & .724 & .319 & .200 & .668 & .305 \\
\arrayrulecolor{kit-gray30}\hline\arrayrulecolor{black}
\multirow{2}{*}{LR} & OpenAI    & .160 & .615 & .253 & .271 & .776 & .402 & .215 & .772 & .334 & .215 & .721 & .330 \\
 & TF-IDF                       & .160 & .523 & .244 & \textBF{.307} & .717 & \textBF{.429} & .221 & .661 & .330 & .229 & .634 & .334 \\
\arrayrulecolor{kit-gray30}\hline\arrayrulecolor{black}
\multirow{2}{*}{NB} & OpenAI    & .164 & .592 & .257 & .272 & .763 & .401 & .252 & .627 & .357 & .230 & .661 & .338 \\
 & SB\textsubscript{max}        & .138 & .549 & .220 & .270 & .699 & .389 & .254 & .683 & .366 & .221 & .644 & .325 \\
\arrayrulecolor{kit-gray30}\hline\arrayrulecolor{black}
\multirow{2}{*}{RF} & OpenAI    & .165 & .588 & .257 & .280 & .722 & .403 & .234 & .746 & .354 & .226 & .685 & .338 \\
 & SB\textsubscript{max}        & .138 & .581 & .222 & .284 & .708 & .405 & .212 & .683 & .322 & .211 & .658 & .316 \\
\arrayrulecolor{kit-gray30}\hline\arrayrulecolor{black}
\multirow{2}{*}{SVM} & OpenAI   & \textBF{.195} & .601 & \textBF{.294} & .280 & \textBF{.803} & .415 & .253 & .726 & .371 & \textbf{.243} & .710 & \textBF{.360} \\
 & SB\textsubscript{max}        & .150 & .591 & .238 & .284 & .733 & .410 & \textBF{.257} & .727 & \textBF{.376} & .230 & .684 & .341 \\
\arrayrulecolor{kit-gray30}\hline\arrayrulecolor{black}
\multirow{2}{*}{XGB} & OpenAI   & .163 & .603 & .257 & .291 & .708 & .412 & .223 & .640 & .329 & .226 & .650 & .333 \\
 & SB\textsubscript{max}        & .139 & .580 & .225 & .279 & .683 & .396 & .222 & .697 & .335 & .213 & .653 & .319 \\
 \arrayrulecolor{kit-gray30}\hline\arrayrulecolor{black}
 \multicolumn{2}{l}{RoBERTa\textsubscript{fine-t.}} & .133 & .593 & .210 & .296 & .656 & .401 & .134 & .822 & .228 & .187 & .690 & .334 \\
\arrayrulecolor{kit-gray30}\hline\arrayrulecolor{black}
\multicolumn{2}{l}{GPT-5-mini}  & .167 & .012 & .023 & .176 & .011 & .020 & .222 & .043 & .073 & .188 & .022 & .066 \\
 \arrayrulecolor{kit-gray30}\hline\arrayrulecolor{black}
 \multicolumn{2}{l}{Random} & .078 & .080 & .078 & .196 & .202 & .199 & .183 & .207 & .191 & .152 & .163 & .156 \\
        \bottomrule
    \end{tabularx}
\end{table}


\ref{RQ2 Performance comparison} examines the performance of classifiers in predicting \acl{ASI} issues.
We evaluated \textbf{supervised machine learning models} (seven classifiers combined with three embeddings), \textbf{finetuned RoBERTa models} and \textbf{Zero-shot LLMs based classifiers}.

For the smell-inducing issues based on \ARCADE (\autoref{tab:rq2_arcade}), \ac{SVM} with OpenAI embeddings achieves the best average performance with an F\textsubscript{1}-score of .360 (P=.243, R=.710).
It exhibits strong recall (.710 average) while maintaining the second highest precision across all three systems, particularly excelling on the StackGres project with an F\textsubscript{1}-score of .415.
Only three configurations achieve higher recall on average (\ac{ffNN}, \ac{kNN}, and \ac{LR}, each with OpenAI embeddings as well).
However, they are not as precise as \ac{SVM}, resulting in lower average F\textsubscript{1}-scores.
% Notably, the \ac{kNN} classifier with OpenAI embeddings achieves the highest recall (.753 average) but suffers from lower precision (.198 average), resulting in an F\textsubscript{1}-score of .309.
% Only the zero-shot classifier using GPT-5-mini achieves slightly higher precision than \ac{SVM} at .248; however, it achieves only 3.9\% recall, failing to identify most smell-inducing issues.
Finetuned RoBERTa models perform relatively similar to supervised ML models. 
Zero-shot classifier using GPT-5-mini performs worse, achieving slightly higher precision than random guessing. 
This demonstrates that the intrinsic architectural reasoning capabilities of \acp{LLM} alone are insufficient without learning from project-specific patterns.
Our significance tests show except zero-shot classifier, all performs significantly better than baseline.
Their p-values can be found in replication package\cite{replication}.

Results for smell-inducing issues based on \Arcan (\autoref{tab:rq2_arcan}) show better overall performance compared to \ARCADE.
\ac{SVM} with OpenAI embeddings remains the best configuration, with an average F\textsubscript{1}-score of .506 (P=.405, R=.741).
This configuration demonstrates strong performance across systems, with the highest individual F\textsubscript{1}-scores on both projects (.400 on StackGres and .612 on Shepard).
% Additional strong performers include \ac{NB} (F\textsubscript{1}=.489) and \ac{RF}, both with OpenAI embeddings (F\textsubscript{1}=.496).
\ac{kNN} with OpenAI embeddings again achieves the highest recall (.782 average) but low precision (.360 average), resulting in an F\textsubscript{1}-score of .475.

Across both datasets, OpenAI embeddings combined with an \ac{SVM} consistently outperform, suggesting that advanced pre-trained embeddings, paired with project-specific training, are most effective for predicting \ac{ASI} issues.

\begin{rqanswer}
    \textbf{Answer to RQ2. }
    \ac{SVM} with OpenAI embeddings achieves the best overall performance, with a balance of precision and recall and the highest average F\textsubscript{1}-scores.
    Zero-shot \acp{LLM} without project-specific training performs poorly.
\end{rqanswer}

\begin{table}
    \centering
    \footnotesize
    \setlength{\tabcolsep}{1.5pt}
    \renewcommand{\arraystretch}{1.2}
    \caption{Results of the Top-2 Configuration of Each Classifier on the Smells Recovered With \Arcan 
    (\acf{ffNN}, \acf{kNN}, \acf{LR}, \acf{NB}, \acf{RF}, \acf{SVM}, \acf{XGB}, and
    %ffNN = feedforward neural network,
    %kNN = k-nearest neighbors,
    %LR = logistic regression,
    %NB = Na\"ive Bayes,
    %RF = random forest,
    %SVM = support vector machine,
    %XGB = XGBoost
    Sentence-BERT (SB)).
    }
    \label{tab:rq2_arcan}
    \begin{tabularx}{\columnwidth}{@{}llZZZZZZZZZ@{}}
        \toprule
                                       &                             & \multicolumn{3}{c}{StackGres} & \multicolumn{3}{c}{Shepard} & \multicolumn{3}{c}{Avg.}                                                                             \\
        \cmidrule(lr){3-5} \cmidrule(lr){6-8} \cmidrule(lr){9-11}
        Class.                     & Emb.                   & \multicolumn{1}{c}{P}                         & \multicolumn{1}{c}{R}                      & \multicolumn{1}{c}{F\textsubscript{1}}       & \multicolumn{1}{c}{P} & \multicolumn{1}{c}{R} & \multicolumn{1}{c}{F\textsubscript{1}} & \multicolumn{1}{c}{P} & \multicolumn{1}{c}{R} & \multicolumn{1}{c}{F\textsubscript{1}} \\
        \midrule
\multirow{2}{*}{ffNN} & OpenAI  & .264 & .692 & .381 & .482 & .713 & .572 & .373 & .702 & .477 \\
 & SB\textsubscript{TF-IDF}     & .247 & .726 & .367 & .452 & .658 & .532 & .349 & .692 & .449 \\
\arrayrulecolor{kit-gray30}\hline\arrayrulecolor{black}
\multirow{2}{*}{kNN} & OpenAI   & .229 & \textBF{.805} & .357 & .490 & \textBF{.758} & .594 & .360 & \textBF{.782} & .475 \\
 & SB\textsubscript{max}        & .219 & .745 & .338 & .470 & .725 & .570 & .345 & .735 & .454 \\
\arrayrulecolor{kit-gray30}\hline\arrayrulecolor{black}
\multirow{2}{*}{LR} & OpenAI    & .255 & .768 & .383 & .501 & .658 & .567 & .378 & .713 & .475 \\
 & TF-IDF                       & .272 & .698 & .390 & .511 & .667 & .576 & .391 & .683 & .483 \\
\arrayrulecolor{kit-gray30}\hline\arrayrulecolor{black}
\multirow{2}{*}{NB} & OpenAI    & .267 & .724 & .389 & .518 & .688 & .589 & .393 & .706 & .489 \\
 & TF-IDF                       & .252 & .798 & .382 & .466 & .654 & .542 & .359 & .726 & .462 \\
\arrayrulecolor{kit-gray30}\hline\arrayrulecolor{black}
\multirow{2}{*}{RF} & OpenAI    & .273 & .730 & .396 & .533 & .683 & .595 & .403 & .707 & .496 \\
 & SB\textsubscript{TF-IDF}     & .265 & .690 & .383 & .456 & .650 & .535 & .360 & .670 & .459 \\
\arrayrulecolor{kit-gray30}\hline\arrayrulecolor{black}
\multirow{2}{*}{SVM} & OpenAI   & .271 & .766 & \textBF{.400} & .538 & .717 & \textBF{.612} & .405 & .741 & \textbf{.506 }\\
 & SB\textsubscript{max}        & .252 & .736 & .375 & .523 & .638 & .570 & .387 & .687 & .472 \\
\arrayrulecolor{kit-gray30}\hline\arrayrulecolor{black}
\multirow{2}{*}{XGB} & OpenAI   & \textBF{.275} & .698 & .394 & .506 & .696 & .583 & .391 & .697 & .488 \\
 & SB\textsubscript{max}        & .254 & .703 & .372 & .501 & .667 & .570 & .377 & .685 & .471 \\
\arrayrulecolor{kit-gray30}\hline\arrayrulecolor{black}
\multicolumn{2}{l}{RoBERTa\textsubscript{fine-t.}} & .243 & .639 & .350 & .451 & .525 & .480 & .347 & .582 & .415  \\
\arrayrulecolor{kit-gray30}\hline\arrayrulecolor{black}
\multicolumn{2}{l}{GPT-5-mini}   & .273 & .013 & .024 & \textbf{.833} & .042 & .080 & \textbf{.553} & .027 & .052 \\
\arrayrulecolor{kit-gray30}\hline\arrayrulecolor{black}
\multicolumn{2}{l}{Random} & .116 & .122 & .118 & ..357 & .346 & .349 & .236 & .234 & .234  \\



        \bottomrule
    \end{tabularx}
\end{table}

\subsection{RQ3: Temporal Information Availability}
\begin{table}
    \centering
    \caption{Avg. Results with Addit. Discussion Input (C\textsubscript{jit}) Compared to Performance on C\textsubscript{0}. Results For the Best Configuration From RQ2 (T\&D) and the Best Performing on the Addit. Data (Full).
    %ffNN = feedforward neural network,
    %kNN = k-nearest neighbors,
    %LR = logistic regression,
    %NB = Na\"ive Bayes,
    %RF = random forest,
    %SVM = support vector machine,
    %XGB = XGBoost;
    %(SD = smell detector) %, TOP specifies whether this configuration was top on full issue (onFull, see RQ2) or on the reduced input (onT\&D). Deltas are calculated based on performance of this configuration in RQ2.
    }
    \setlength{\tabcolsep}{1.5pt}
    \renewcommand{\arraystretch}{1.1}
    \label{tab:rq3}
    \begin{tabularx}{\columnwidth}{lZZXcccccc}
        \toprule
                                     &&  &           &                \multicolumn{6}{c}{Average}                                                                                                           \\
        \cmidrule(lr){5-10}
        \multirow{-2}{3.2em}{Smell- detector}&\multirow{-2}{3.2em}{Config. top on}&Class.                     & Emb.                & \multicolumn{1}{c}{P} & $\Delta$                       & \multicolumn{1}{c}{R}     & $\Delta$                   & \multicolumn{1}{c}{F\textsubscript{1}}& $\Delta$          \\
        \midrule
        \multirow{2}{*}{\ARCADE}&T\&D&\multirow{2}{*}{SVM} & \multirow{2}{*}{OpenAI}& \multirow{2}{*}{.241}	&\multirow{2}{*}{(\textcolor{kit-red}{--.002})}&\multirow{2}{*}{.721}	&\multirow{2}{*}{(\textcolor{kit-green}{+.011})}&\multirow{2}{*}{.360}&\multirow{2}{*}{($\pm$.000)}
\\
        &Full&\\
         \arrayrulecolor{kit-gray30}\midrule\arrayrulecolor{black}
        \multirow{2}{*}{\Arcan}&T\&D&\multirow{2}{*}{SVM} & \multirow{2}{*}{OpenAI}& \multirow{2}{*}{.406}	&\multirow{2}{*}{(\textcolor{kit-green}{+.001})}&\multirow{2}{*}{.737}	&\multirow{2}{*}{(\textcolor{kit-red}{--.004})}&\multirow{2}{*}{.505}&\multirow{2}{*}{(\textcolor{kit-red}{$-$.001})}
\\
        &Full&\\
        \bottomrule
    \end{tabularx}
\end{table}

\ref{RQ3 Tempral availiabilty} investigates the impact of information availability on prediction performance by comparing two temporal conditions: C\textsubscript{0} (title and description only, available at issue creation, as used in \ref{RQ2 Performance comparison}) versus C\textsubscript{jit} (including pre-implementation comments).
\autoref{tab:rq3} presents results for the best-performing configuration (SVM with OpenAI embeddings).

The results show minimal differences between the two conditions.
For smell-inducing issues based on \ARCADE detection, additionally using pre-implementation discussions (C\textsubscript{jit}) achieves nearly identical performance to using only title and description (C\textsubscript{0}): precision changes by +.002 to .241, recall increases by +.011 to .721, and F\textsubscript{1}-score remains at .360.
For \Arcan, the added information yields a marginal decrease in average F\textsubscript{1}-score to .505 (compared to .506), with precision at .406 (+.001) and recall at .737 (-.004).

%These minimal performance differences suggest that issue titles and descriptions alone contain sufficient information for \ac{ASI} issues and that adding pre-implementation discussions contributes negligibly.
%This finding has practical implications: predictions can be made immediately upon issue creation, without waiting for community discussion to accumulate, enabling earlier architectural awareness.
%The consistency of this pattern across both smell detectors strengthens the generalizability of this observation.

\begin{rqanswer}
    \textbf{Answer to RQ3.}
    Issue titles and descriptions alone contain sufficient information for predicting \ac{ASI} issues; adding pre-implementation discussions contributes negligibly.
\end{rqanswer}

\subsection{RQ4: Cross-Project Generalizability}
% https://docs.google.com/spreadsheets/d/1iTPQteP3PXsS0o2DIOu_qRHoheUE8BjS2B7A8Cx0Sjg/edit?usp=sharing

\begin{table}
  \centering
  \caption{Generalization Performance in Cross-Project Setting (Relative Change to RQ2)}
  \label{tab:rq4}
  \setlength{\tabcolsep}{4pt} % compact columns
  \renewcommand{\arraystretch}{1} % compact rows
    \begin{tabular}{llccc}
      \toprule
      Detector & Dataset & Precision ($\Delta$) & Recall ($\Delta$) & F\textsubscript{1} ($\Delta$) \\
      \midrule
      \multirow{4}{*}{\ARCADE}
        & Inkscape  & .115 (\textcolor{kit-red}{$-.080$}) & .070 (\textcolor{kit-red}{$-.531$}) & .087 (\textcolor{kit-red}{$-.207$}) \\
        & StackGres & .209 (\textcolor{kit-red}{$-.071$}) & .297 (\textcolor{kit-red}{$-.506$}) & .246 (\textcolor{kit-red}{$-.169$}) \\
        & Shepard   & .181 (\textcolor{kit-red}{$-.072$}) & .739 (\textcolor{kit-green}{$+.013$}) & .291 (\textcolor{kit-red}{$-.080$}) \\
        \arrayrulecolor{kit-gray30}\cmidrule(){2-5}\arrayrulecolor{black}
        & Average   & .168 (\textcolor{kit-red}{$-.074$}) & .369 (\textcolor{kit-red}{$-.341$}) & .208 (\textcolor{kit-red}{$-.152$}) \\
       
               \arrayrulecolor{kit-gray30}\midrule\arrayrulecolor{black}
      \multirow{3}{*}{\Arcan}
       & StackGres & .210 (\textcolor{kit-red}{$-.061$}) & .561 (\textcolor{kit-red}{$-.205$}) & .305 (\textcolor{kit-red}{$-.095$}) \\
        & Shepard   & .439 (\textcolor{kit-red}{$-.099$}) & .150 (\textcolor{kit-red}{$-.567$}) & .224 (\textcolor{kit-red}{$-.388$}) \\
        \arrayrulecolor{kit-gray30}\cmidrule(){2-5}\arrayrulecolor{black}
        & Average & .324 (\textcolor{kit-red}{$-.080$}) & .356 (\textcolor{kit-red}{$-.386$}) & .265 (\textcolor{kit-red}{$-.242$}) \\
       
      \bottomrule
    \end{tabular}%

\end{table}
\ref{RQ4 Cross project} investigates the transferability of trained models across different projects.
\autoref{tab:rq4} presents results for the best-performing configuration from RQ2 (SVM with OpenAI embeddings) in cross-project settings.
The results show substantial performance degradation in cross-project scenarios compared to within-project evaluation, although still better than random baseline.
%
For \ARCADE-based smells, all three projects show reduced performance.
%Inkscape shows F\textsubscript{1}-score of .087, StackGres achieves .246, and Shepard reaches .291.
Notably, Shepard's recall increases by 2\% to .739, maintaining the ability to identify \ac{ASI} issues, but suffers a 40\% precision loss (to .181) and a 27\% decline in F\textsubscript{1}-score.
%
For \Arcan-based smells, both projects show performance losses.
%These performance losses across all projects suggest that textual patterns indicative of \acl{ASI} issues are highly project-specific.
%The lack of generalization indicates that architectural vocabulary, development practices, and issue description styles vary between projects, limiting the applicability of models trained on external codebases. 
\begin{rqanswer}
    \textbf{Answer to RQ4:} Cross-project models generalize poorly and only slightly outperform the baseline.
\end{rqanswer}



